
\todo{Future work: extend the comparison with \acrfull{fde} and the proposed novel clear-box method/metric, kept as a thread to pick up later.}
As \acrfull{ai} systems move from controlled laboratory settings into everyday and safety-critical applications, the ability to tell when a model is likely to be wrong becomes as important as its raw accuracy. Transformer-based \acrfull{asr} models such as Whisper achieve low error rates, yet they offer no native signal of how reliable each individual transcription is, and they are known to fail silently, including fluent but fabricated output. This work studies \acrfull{uq} for Whisper on Spanish speech, where training data is comparatively scarce.

We compare three post-hoc \acrshort{uq} methods: \acrfull{ts}, \acrfull{mcd}, and a Levenshtein-based variant we refer to as \acrfull{lmcd}. Each method produces a per-utterance uncertainty score, and we measure its quality as the Pearson correlation between that score and the \acrfull{wer} of the transcription. Because this objective is non-differentiable, noisy, and expensive to sample, every method is tuned on equal footing with \acrfull{cmaes}, a derivative-free optimizer, before the methods are compared.

Across ten folds of a Spanish \acrshort{uq} benchmark built from Mexican and Chilean speech, \acrshort{lmcd} produces uncertainty scores that track transcription error substantially better than \acrshort{ts} and \acrshort{mcd}, and is the only method that separates from the others by a statistically significant margin. Tuning also reveals that the dropout-based methods only become informative at dropout rates well below the usual defaults, and that \acrshort{lmcd} and \acrshort{mcd} stay stable between calibration and test data whereas \acrshort{ts} degrades the most. These results position distance-based, sample-agreement scoring as a practical basis for flagging likely transcription errors.

